% ============================================================
%  Assignment 3 — Regularization and Stability
%  Statistical Methods for Machine Learning — A.Y. 2025/26
%  Instructor: Nicolò Cesa-Bianchi
% ============================================================
\documentclass[11pt, a4paper, twoside]{article}
\raggedbottom

\usepackage[T1]{fontenc}
\usepackage[utf8]{inputenc}
\usepackage[english]{babel}


\usepackage[margin=2.2cm, top=2.5cm, bottom=2.5cm]{geometry}
\usepackage{amsmath, amssymb, amsthm}
\usepackage{mathtools}
\usepackage{graphicx}
\usepackage{float}
\usepackage{booktabs}
\usepackage{xcolor}
\usepackage{hyperref}
\usepackage{cleveref}
\usepackage{titlesec}

% ── section style: compact, like JMLR ────────────────────────
\titleformat{\section}{\normalfont\large\bfseries}{\thesection.}{0.5em}{}
\titleformat{\subsection}{\normalfont\normalsize\bfseries}{\thesubsection.}{0.5em}{}
\titleformat{\paragraph}[runin]{\normalfont\normalsize\bfseries}{}{0pt}{}[.]
\titlespacing{\section}{0pt}{2.4ex plus .4ex}{0.8ex plus .1ex}
\titlespacing{\subsection}{0pt}{1.8ex plus .3ex}{0.6ex plus .1ex}
\titlespacing{\paragraph}{0pt}{0.8ex plus .1ex}{0.5em}

\hypersetup{colorlinks=true, linkcolor=blue!60!black,
            citecolor=green!50!black, urlcolor=blue!70!black}

% ── theorem environments ─────────────────────────────────────
\newtheorem{definition}{Definition}
\newtheorem{theorem}{Theorem}
\newtheorem{remark}{Remark}

% ── math macros ──────────────────────────────────────────────
\newcommand{\R}{\mathbb{R}}
\newcommand{\E}{\mathbb{E}}
\newcommand{\norm}[1]{\left\lVert #1 \right\rVert}
\newcommand{\abs}[1]{\left| #1 \right|}
\newcommand{\AS}{A_S}
\newcommand{\ASi}{A_{S^{\setminus i}}}
\newcommand{\betahat}{\hat{\beta}}

% ── title ────────────────────────────────────────────────────
\title{%
    \vspace{-1em}%
    \textbf{Assignment 3 --- Regularization and Stability}\\[3pt]
    \normalsize Statistical Methods for Machine Learning --- A.Y.\ 2025/26\\
    \normalsize Instructor: Nicolò Cesa-Bianchi
}
\author{%
    Lize Ana Zabote\\
    \small Student ID: V13380 \quad
}
\date{}

% ============================================================
\begin{document}
% ============================================================

\maketitle
\vspace{-1em}

\noindent\textbf{Declaration.}
\textit{I declare that this material, which I now submit for assessment, is entirely my own
work and has not been taken from the work of others, save and to the extent that such work
has been cited and acknowledged within the text of my work. I understand that plagiarism,
collusion, and copying are grave and serious offences in the university and accept the
penalties that would be imposed should I engage in plagiarism, collusion or copying. This
assignment, or any part of it, has not been previously submitted by me or any other person
for assessment on this or any other course of study.}

\medskip

\begin{abstract}
\noindent
This report looks at how regularization, algorithmic stability, and generalization
interact in linear regression, using the framework of Bousquet and
Elisseeff~\cite{bousquet2002} as a guide.
Ridge regression and OLS are implemented from scratch and stability is measured through a leave-one-out
retraining procedure applied to both synthetic dataset ($n=300$, $d=20$) and the
Diabetes benchmark.
The results match the theory: Ridge stability $\betahat$ decreases as
$\lambda$ grows, in line with the $\beta = \mathcal{O}(1/\lambda m)$ bound.
On the Diabetes data, Ridge finds a clear optimum around $\lambda^* \approx 133$
and outperforms OLS. On the synthetic data, OLS actually does slightly better at
moderate $\lambda$.
A secondary experiment varying training set size shows that $n$ and $\lambda$ play
symmetric roles in the stability bound.
\end{abstract}

\medskip\noindent\rule{\linewidth}{0.4pt}\medskip

% ============================================================
\section{Introduction}
% ============================================================

A recurring question in machine learning is why regularized algorithms generalize
well even when the training set is small or noisy.
The standard answer from VC theory~\cite{vapnik1982} bounds generalization error
through the complexity of the hypothesis class, but this approach runs into trouble
for algorithms like regularized regression or $k$-nearest neighbors, where the search
space is not easily characterized by a single complexity measure.

Bousquet and Elisseeff~\cite{bousquet2002} take a different angle.
Rather than asking how large the model class is, they ask how much the learned
function changes when one training point is removed.
If the change is minimum, the algorithm is stable, and stability turns out to
imply a small gap between training and generalization error.
What makes their result particularly useful for this assignment is the bound in
Section~5.2 of~\cite{bousquet2002}: for regularization-based algorithms,
$\beta = \mathcal{O}(1/\lambda m)$.
It gives a direct, quantitative link between $\lambda$, sample size, and stability
that can actually be tested.

Ridge regression and OLS are implemented from scratch using only \texttt{numpy},
stability is estimated by retraining models after removing each training point in
turn, and the results are compared across datasets and regularization strengths.

% ============================================================
\section{Theoretical Background}
% ============================================================

\paragraph{Setup}
Let $S = \{z_1, \ldots, z_m\}$ with $z_i = (x_i, y_i)$ drawn i.i.d.\ from $\mathcal{D}$.
A learning algorithm $A$ maps $S$ to a function $\AS : \mathcal{X} \to \mathcal{Y}$.
Given a loss $\ell(f, z) = c(f(x), y)$, the three quantities used throughout are:
\begin{align*}
    R \;=\; \E_z[\ell(\AS, z)], \quad
    R_\text{emp} \;=\; \tfrac{1}{m}\textstyle\sum_i \ell(\AS, z_i), \quad
    R_\text{loo} \;=\; \tfrac{1}{m}\textstyle\sum_i \ell(\ASi, z_i),
\end{align*}
where $S^{\setminus i}$ is the training set with the $i$-th point removed.

\paragraph{Notions of stability}
The paper introduces three nested stability notions.
Hypothesis stability (Definition~3) bounds how much the average loss changes when one
training point is removed; error stability (Definition~5) bounds the change in
expected risk; uniform stability (Definition~6) is the strongest --- it bounds the
\emph{worst-case} pointwise change in loss over all possible test points:

\begin{definition}[Uniform stability,~\cite{bousquet2002}]
Algorithm $A$ has \emph{uniform stability} $\beta$ with respect to $\ell$ if
\[
    \forall S \in \mathcal{Z}^m,\; \forall i \in \{1,\ldots,m\},\quad
    \sup_{z \in \mathcal{Z}} \abs{\ell(\AS, z) - \ell(\ASi, z)} \leq \beta.
\]
\end{definition}

Uniform stability implies hypothesis stability, which implies error stability.
An algorithm is called \emph{stable} when $\beta_m = \mathcal{O}(1/m)$.

\paragraph{Generalization bound}

\begin{theorem}[Theorem~12 in~\cite{bousquet2002}]
Let $A$ have uniform stability $\beta$ with respect to a loss $\ell$ such that
$0 \leq \ell(\AS, z) \leq M$ for all $z \in \mathcal{Z}$ and all $S$.
For any $m \geq 1$ and $\delta \in (0,1)$, with probability at least $1-\delta$
over the random draw of $S$:
\[
    R \;\leq\; R_\text{emp} + 2\beta +
    (4m\beta + M)\sqrt{\tfrac{\ln(1/\delta)}{2m}}.
\]
\end{theorem}

The bound controls $R - R_\text{emp}$, the gap between
generalization and training error, rather than $R$ itself.
This distinction turns out to matter quite a bit in practice, as the experiments in
Section~\ref{sec:discussion} show.

\paragraph{Stability of Ridge regression}
For regularized least squares in a reproducing kernel Hilbert space (RKHS) with bounded
labels $Y = [0, B]$ and kernel satisfying $k(x,x) \leq \kappa^2$, Example~3
of~\cite{bousquet2002} gives the uniform stability bound:
\[
    \beta \;\leq\; \frac{2\kappa^2 B^2}{\lambda m},
\]
The bound scales as $\mathcal{O}(1/\lambda m)$: larger $\lambda$ or more training
data both reduce $\beta$. This is the central quantitative prediction verified
empirically in Section \textbf{5}.

% ============================================================
\section{Datasets}
% ============================================================

\paragraph{Synthetic}
The synthetic dataset consists of $n = 300$ samples in $d = 20$ dimensions,
generated as:
$w^* \sim \mathcal{N}(0, I_d)$,\;
$x_i \sim \mathcal{N}(0, I_d)$,\;
$y_i = x_i^\top w^* + \varepsilon_i$,\; $\varepsilon_i \sim \mathcal{N}(0, 2^2)$.
The ratio $n/d = 15$ keeps the problem well in the underparameterized regime, so
$X^\top X$ is well-conditioned and OLS is numerically stable.
The noise level $\sigma = 2$ is deliberately high to make the bias-variance tradeoff
visible: without some noise, any reasonable estimator would fit well and differences
between Ridge and OLS would be hard to see.
The features and target are standardized on the training set only and a 75/25 split is
applied with a fixed seed.

\paragraph{Diabetes (real-world)}
The Diabetes dataset (442 samples, 10 features) from \texttt{sklearn.datasets} serves
as a real-world benchmark~\cite{efron2004}.
Compared to the synthetic data, its features are moderately correlated --- which is precisely the scenario where Ridge regularization actually helps.
The same preprocessing protocol is applied.

% ============================================================
\section{Methods}
% ============================================================

\paragraph{Algorithms}
All algorithms are implemented using \texttt{numpy}; no \texttt{sklearn}
estimator \texttt{.fit()} calls are used.
\emph{Ridge regression} uses the closed-form solution
$w_\lambda = (X^\top X + \lambda I_d)^{-1} X^\top y$,
solved via \texttt{numpy.\allowbreak linalg.\allowbreak solve} for numerical stability.
\emph{OLS} uses Moore-Penrose pseudoinverse
(\texttt{numpy.\allowbreak linalg.\allowbreak lstsq},
equivalent to Ridge with $\lambda \to 0$) as the unregularized baseline.

\paragraph{Stability estimation}
Stability is estimated empirically using the following leave-one-out metric:
\begin{equation}
    \betahat \;=\; \frac{1}{m} \sum_{i=1}^{m}
    \frac{1}{|S_\text{test}|}
    \sum_{x \in S_\text{test}} \abs{f_S(x) - f_{S^{\setminus i}}(x)},
    \label{eq:beta}
\end{equation}
where $f_S = \AS$ is the model trained on the full training set and
$f_{S^{\setminus i}} = \ASi$ is retrained after removing point $i$.
For each training point, the model is retrained without it, and the average
absolute shift in predictions over the test set is recorded. The metric $\betahat$
is the mean of these shifts across all $m$ leave-one-out steps.

In the theory, uniform stability (Definition~6) takes the
\emph{supremum} of the loss change over all test points, while hypothesis stability
(Definition~3) averages it. The metric in Eq.~\eqref{eq:beta} averages the change in
\emph{predictions} rather than losses, so it is somewhere between the two metrics.
Since the squared loss is Lipschitz and both datasets are standardized, the two
quantities are proportional, and $\betahat$ works as a reasonable proxy for $\beta$.

\paragraph{OLS stability}
OLS has no regularization parameter, so $\betahat$ does not vary with $\lambda$ for
this model. It appears in the plots as a flat reference line, representing the
baseline sensitivity of the unregularized estimator.

\paragraph{Implementation}
The project is organized as follows:
\texttt{src/models.py} implements Ridge and OLS from scratch;
\texttt{src/datasets.py} loads and preprocesses both datasets;
\texttt{src/stability.py} implements Eq.~\eqref{eq:beta} and the experiment runners;
\texttt{src/plotting.py} produces one PDF figure per experiment;
\texttt{main.py} is the single entry point that runs all experiments.
Results are fully reproducible by running \texttt{python main.py}
with a fixed \texttt{numpy.random.seed(42)}.

% ============================================================
\section{Results}
% ============================================================

\subsection{Stability vs.\ $\lambda$}

Figures~\ref{fig:synth_stab} and~\ref{fig:diab_stab} show $\betahat$ as a function
of $\lambda \in [10^{-4}, 10^3]$ for both datasets.

\begin{figure}[H]
    \centering
    \includegraphics[width=0.78\linewidth]{fig1_synthetic_stability.pdf}
    \caption{Synthetic dataset --- Stability $\betahat$ vs.\ $\lambda$.
             Ridge $\betahat$ decreases monotonically from approximately $0.00667$
             at $\lambda = 10^{-4}$ to $0.00198$ at $\lambda = 10^3$,
             consistent with the $\mathcal{O}(1/\lambda m)$ bound.
             }
    \label{fig:synth_stab}
\end{figure}

\begin{figure}[H]
    \centering
    \includegraphics[width=0.78\linewidth]{fig3_diabetes_stability.pdf}
    \caption{Diabetes dataset --- Stability $\betahat$ vs.\ $\lambda$.
             The same qualitative pattern holds: Ridge $\betahat$ decreases
             monotonically from approximately $0.00390$ to $0.00124$ as $\lambda$
             increases.}
    \label{fig:diab_stab}
\end{figure}

On both datasets, Ridge $\betahat$ decreases monotonically with $\lambda$,
which is the direct empirical counterpart of $\beta \leq 2\kappa^2 B^2 / (\lambda m)$.
At small $\lambda$, Ridge and OLS have nearly identical stability,
since Ridge approaches OLS as $\lambda \to 0$.
As $\lambda$ grows, the Ridge curve drops below the OLS reference line, showing
that regularization does produce a more stable estimator.

\subsection{Train and test error vs.\ $\lambda$}

\begin{figure}[H]
    \centering
    \includegraphics[width=0.78\linewidth]{fig2_synthetic_error.pdf}
    \caption{Synthetic dataset --- Train (dashed) and test (solid) MSE vs.\ $\lambda$.
             In well-conditioned underparameterized regime ($n=300$, $d=20$),
             OLS achieves lower test MSE ($\approx 0.31$) than Ridge at any
             moderate $\lambda$, and the Ridge test error rises for
             $\lambda \gtrsim 10$. When $n \gg d$ and the design matrix is well-conditioned, OLS is
             already near its optimal bias-variance tradeoff.}
    \label{fig:synth_err}
\end{figure}

\begin{figure}[H]
    \centering
    \includegraphics[width=0.78\linewidth]{fig4_diabetes_error.pdf}
    \caption{Diabetes dataset --- Train (dashed) and test (solid) MSE vs.\ $\lambda$.
             A clear U-shaped test-error curve is visible, with minimum test MSE
             $\approx 0.492$ at $\lambda^* \approx 133$, compared to OLS test MSE
             $\approx 0.509$. The training MSE rises monotonically with $\lambda$
             as the model is increasingly penalized.}
    \label{fig:diab_err}
\end{figure}

The two datasets present different results.
On the synthetic data, OLS achieves lower test MSE than Ridge for the whole range
of moderate $\lambda$; Ridge only catches up at very large $\lambda$, by which point
both models are already heavily underfitted.
On Diabetes dataset, Ridge outperforms OLS at $\lambda^*$,
where correlated features and a lower signal-to-noise ratio make the variance
reduction from regularization worth the bias it introduces.

Table~\ref{tab:summary} makes the stability-vs-error tradeoff concrete for the
Diabetes dataset. The row at $\lambda \approx 133$ is the one that minimizes test
MSE; the row at $\lambda = 10^3$ has the best stability but the worst predictions.

\begin{table}[H]
\centering
\caption{Ridge stability $\betahat$ and test MSE at selected $\lambda$ values
         (Diabetes dataset).}
\label{tab:summary}
\vspace{4pt}
\begin{tabular}{lcc}
\toprule
$\lambda$         & Stability $\betahat$ & Ridge Test MSE \\
\midrule
$10^{-4}$         & 0.00390              & 0.509          \\
$10^{0}$          & 0.00380              & 0.510          \\
$\approx 133$     & 0.00168              & \textbf{0.492} \\
$10^{3}$          & 0.00124              & 0.559          \\
\bottomrule
\end{tabular}
\end{table}

Going from $\lambda = 10^{-4}$ to $\lambda = 10^3$ roughly triples stability but
increases test MSE by about 10\%.
The reason is that Theorem~12 bounds $R - R_\text{emp}$, not $R$. In other words, a tight
train-test gap around a bad model is still a bad model.

\subsection{Extension: effect of training set size on stability}

Figure~\ref{fig:size} shows $\betahat$ vs.\ $\lambda$ for Ridge regression on the
synthetic dataset ($d=10$) at three training set sizes $n \in \{50, 150, 300\}$.

\begin{figure}[H]
    \centering
    \includegraphics[width=0.78\linewidth]{fig5_size_study.pdf}
    \caption{Stability vs.\ $\lambda$ for Ridge at three training set sizes
             ($d = 10$, synthetic data).
             For any fixed $\lambda$, $\betahat$ is smaller for larger $n$,
             consistent with $\beta \propto 1/({\lambda m})$.
             The three curves converge at large $\lambda$ (regularization dominates)
             and diverge at small $\lambda$.}
    \label{fig:size}
\end{figure}

The curves confirm both dimensions of the bound.
Fixing $\lambda$ and increasing $n$ shifts the curve down; fixing $n$ and increasing
$\lambda$ moves along the curve to the right and down.
The two levers are interchangeable from a stability standpoint: the same $\betahat$
target can be hit either by collecting more data or by regularizing more heavily.

% ============================================================
\section{Discussion}
\label{sec:discussion}
% ============================================================

\paragraph{Regularization increases stability}
On both datasets, $\betahat$ drops as $\lambda$ grows, matching the bound
$\beta \leq 2\kappa^2 B^2/(\lambda m)$ from Example~3 of~\cite{bousquet2002}.
The behaviour near $\lambda = 0$ is also consistent. As Ridge approaches OLS, the
two stability curves converge.
The gap between them is a direct measure of what regularization buys in terms of
reduced sensitivity to individual training points.
Without the $\lambda I$ term, the pseudoinverse can shift noticeably when any single
data point changes the conditioning of $X^\top X$, and this is captured in the
consistently higher $\betahat$ for OLS.

\paragraph{Stability controls the gap, not the absolute error}
This is the most practically relevant point from~\cite{bousquet2002} and
something the Diabetes results illustrate well.
At $\lambda = 10^3$, the model achieves its lowest $\betahat = 0.00124$ --- meaning
train and test error are very close together --- but the test MSE is 0.559, the worst
in the table.
The model has been shrunk so hard toward zero that it fits neither the training nor
the test data well.
The optimal $\lambda^* \approx 133$ is where bias and variance balance out; beyond
that point, adding more regularization improves stability but makes the
model less useful in practice.

\paragraph{When OLS generalizes as well as Ridge}
On the synthetic dataset ($n=300$, $d=20$), OLS actually outperforms Ridge at moderate
$\lambda$.
In this regime $n \gg d$, so $X^\top X$ is far from singular, no individual training
point has unusual leverage, and the pseudoinverse already gives a near-optimal fit.
Adding $\lambda I$ introduces bias without producing any meaningful variance reduction,
so Ridge ends up worse.
This is not a contradiction with~\cite{bousquet2002}: their bound guarantees a
\emph{small train-test gap} for a stable algorithm, not a small test error.
The benefit of regularization only shows up when features are correlated or when $n$
is closer to $d$, as in the Diabetes case.

\paragraph{Dataset size and stability}
The size study confirms $\betahat \propto 1/(\lambda m)$.
One thing worth noting is that the two levers are not equivalent from a modelling
perspective: collecting more data improves stability \emph{and} the quality of the
fit, while increasing $\lambda$ past $\lambda^*$ improves stability at the direct
cost of prediction quality.
In practice this means that regularization is not a substitute for data, even if
it looks like one in the bound.

\paragraph{Limitations}
The leave-one-out procedure requires $m$ retraining steps per $\lambda$ value which kept the dataset sizes
small ($n \leq 300$) to stay within reasonable runtime.
Scaling to larger problems would need some approximation, such as subsampling which
training points to remove.
The metric itself approximates hypothesis stability rather than uniform stability,
since it averages over the test set rather than taking a supremum.

% ============================================================
\section{Conclusion}
% ============================================================

The experiments support the main predictions of Bousquet and
Elisseeff~\cite{bousquet2002}.
Ridge stability $\betahat$ drops monotonically with $\lambda$ on both datasets,
consistent with $\beta = \mathcal{O}(1/\lambda m)$, and the size study shows the
same effect from increasing $n$, confirming that the two parameters enter the bound
symmetrically.
On Diabetes, the best test error is achieved at an intermediate $\lambda^*$, not at
maximum stability, which illustrates that the bound controls the
train-test gap, not the absolute risk.
On the synthetic data, OLS matches or beats Ridge at moderate $\lambda$, which is
the expected outcome when $n \gg d$ and the design matrix is well-conditioned.
Taken together, the results show that stability is a useful lens for understanding
regularization, but it needs to be interpreted alongside the actual test error:
a tight generalization gap only translates to good predictions when the model is
fitting the right signal in the first place.

% ============================================================
\bibliographystyle{plain}
\begin{thebibliography}{9}

\bibitem{bousquet2002}
O.~Bousquet and A.~Elisseeff.
\newblock Stability and generalization.
\newblock \emph{Journal of Machine Learning Research}, 2:499--526, 2002.

\bibitem{vapnik1982}
V.~N. Vapnik.
\newblock \emph{Estimation of Dependences Based on Empirical Data}.
\newblock Springer-Verlag, 1982.

\bibitem{efron2004}
B.~Efron, T.~Hastie, I.~Johnstone, and R.~Tibshirani.
\newblock Least angle regression.
\newblock \emph{The Annals of Statistics}, 32(2):407--499, 2004.

\end{thebibliography}

\end{document}